\documentclass[11pt]{article}

\usepackage[T1]{fontenc}
\usepackage{lmodern}
\usepackage{amsmath}
\usepackage{amssymb}
\usepackage[margin=1in]{geometry}
\usepackage{booktabs}
\usepackage{hyperref}

\newcommand{\key}[1]{\texttt{#1}}
\newcommand{\dd}{\mathrm{d}}

\title{Electron-scattering acceleration by a modified random walk:\\
the step as implemented, in Cartesian and spherical-polar cells, at rest, in a moving
medium and in curved spacetime}
\author{Athena++ Monte Carlo radiation transport}
\date{}

\begin{document}
\sloppy
\maketitle

\section{What the step does and where it lives}

A photon deep inside an optically thick cell scatters many times before it leaves.
The modified random walk (MRW) of Fleck \& Canfield (1984), as used by Min et al.\
(2009) and simplified by Robitaille (2010), replaces those scatterings by one draw from
the diffusion equation: the time the photon takes to reach the boundary of a region in
which the medium is uniform, and the point on that boundary where it arrives.  For
Thomson scattering that is the whole step.  For Compton scattering the photon's energy
also changes along the way; given the Thomson path $\tau_{\rm path}$ the walk covered,
the final energy is drawn from the Green's function of the Kompaneets equation at
Compton parameter $y=\theta\tau_{\rm path}$, $\theta=kT/m_ec^2$, and with free-free
absorption the photon carries the survival fraction $S(y)$ as a weight.  The Green's
function, its absorbed form, the quantile tables and their validation are the subject
of \texttt{doc/monte\-\_carlo/kompaneets/kompaneets\_greens\_method.tex} and are not
repeated here; this document is about the step itself.

The code is in \texttt{src/monte\_carlo/mrw.cpp} (the step and its tables), with the
table reader and sampler of the Green's function in \texttt{kgreens.cpp}.  The pushers
call it from the branch of their transport loop that would otherwise place an
interaction inside the current cell: Cartesian and spherical-polar pushers through
\texttt{MRWTrigger} and \texttt{MRWStep}, the general pusher through the same two with
its own frame construction (\texttt{GeneralPusher::MRWFrame}).  The step ends the free
flight: on return the photon has a new position, an isotropic new direction, an energy
and a weight, and the pusher draws a fresh optical depth and continues, so the block
loop never treats the step as an interaction.  The input keys are collected in
Section~\ref{sec:keys}.

Everything in the step is dimensionless until the last moment, so one set of tables
serves every cell of every run.  For a diffusion coefficient $D=1/(3\chi)$, with $\chi$
the comoving extinction, and a domain of size $L$, the natural time is
\begin{equation}
  s=\frac{\pi^2 D\,ct}{L^2},
\end{equation}
in which the slowest mode of the absorbing-boundary problem decays as $e^{-s}$ and the
Thomson path is $\tau_{\rm path}=\chi\,ct$.

\section{A fixed-time step rather than a fixed-radius one}
\label{sec:fixedtime}

The textbook step draws the first-passage time to a sphere of fixed radius $R_0$ and
places the photon on it.  That is the right step only when the medium is at rest on the
grid.  Three things want the time chosen first and the position second:
\begin{itemize}
\item In a moving medium the fluid element advects by $\beta\gamma ct$ in the lab over
  comoving time $t$, and for a fast enough flow the region reaches a cell face by
  advection before the photon reaches its surface by diffusion.
\item In a dynamic run the photon has a time budget, the remainder of the cycle
  (\texttt{Photon::dtp}), which a first-passage draw ignores.
\item In curved spacetime the comoving frame is defined at an instant and the cell is
  a patch of spacetime; a time is the natural thing to bound.
\end{itemize}
So the step is: choose a comoving time budget $t_b$; with probability $1-P(t_b)$, where
$P$ is the fraction of walkers still inside the domain at $t_b$, the photon reached the
boundary earlier, at a time $t_e$ drawn from the first-passage distribution conditioned
on $t_e<t_b$, and is placed on the boundary; otherwise it is still inside at $t_b$, at a
position drawn from the fixed-time solution, and $t_e=t_b$.  Both outcomes are isotropic
in the new direction.  With no budget and a medium at rest the first branch is always
taken and the step reduces to the fixed-radius one of the references.  The budget is the
smallest of the advection room (Section~\ref{sec:advection}), the cycle remainder, and
the time at which $P$ has fallen to \key{accel\_pmax} (0 by default, meaning no cap).

\section{The domain: the cell itself, or a sphere inside it}
\label{sec:domain}

\subsection{Why the sphere leaves a dead zone}

The sphere of the references is the largest that fits around the photon inside its
cell, so a step is possible only where the distance to the nearest face exceeds the
threshold, a fraction $((L-\tau_a)/L)^3$ of a cubic cell of optical half-width $L$ and
threshold $\tau_a$; and every step lands the photon on the sphere, at distance zero from
a face, from where it random-walks about $\tau_a^2$ scatterings before it can step again.
On the gridded sphere of the tests ($8^3$ cells) the sphere domain replaced 13 percent of
the scatterings at $\tau=100$ ($L=15.6$, $\tau_a=10$) and 66 percent at $\tau=1000$
($L=156$, $\tau_a=20$), each within a few percent of the active-volume estimate.  The
sphere is kept as \key{accel\_domain = sphere} for comparison and is not the default.

\subsection{The cell as a box}

For a rectangular box with absorbing walls the diffusion Green's function separates into
three one-dimensional problems, so a photon anywhere in a thick cell can step, and one
near a face steps out through it.  In one dimension, with walls at $0$ and $W$, start
$\xi=x_0/W$ and time $\sigma=\pi^2Dct/W^2$:
\begin{align}
  S(\xi,\sigma)&=\sum_{n\ \rm odd}\frac{4}{n\pi}\sin(n\pi\xi)\,e^{-n^2\sigma}
  &&\text{survival},\\
  g(\xi,\sigma)&=-\frac{\partial\ln S}{\partial\sigma}
  &&\text{hazard},\\
  p_0(\xi,\sigma)&=\frac{\sum_n n\sin(n\pi\xi)e^{-n^2\sigma}}
                       {2\sum_{n\ \rm odd}n\sin(n\pi\xi)e^{-n^2\sigma}}
  &&\text{probability that the exit is through the wall at 0},\\
  C(x\,|\,\xi,\sigma)&=\frac1S\sum_n\frac{2}{n\pi}\sin(n\pi\xi)\,[1-\cos(n\pi x/W)]\,e^{-n^2\sigma}
  &&\text{position CDF, walker still inside}.
\end{align}
The survival of the box is the product of the three one-dimensional survivals, each with
its own $\sigma_i=\pi^2Dct/W_i^2$, so the first-passage time of the box is the solution
of $\prod_iS_i(\xi_i,\sigma_i(t))=1-u$ for a uniform $u$; the dimension through which the
photon leaves is drawn in proportion to the hazard rates $\sigma_i'\,g_i$, the side from
$p_0$, and the positions in the other two dimensions from $C$ at the exit time.  A photon
still inside at the budget gets all three positions from $C$.

All four functions are tabulated once per rank on a grid in $\xi\in[1/130,1/2]$ (the
rest by the symmetry $\xi\to1-\xi$) and $\ln\sigma\in[\ln10^{-4},\ln40]$, the position
quantiles from $\sigma=0.01$ up at 65 levels in single precision.  Outside the grid:
below $\sigma_{\min}$ the two-image solution in $\mathrm{erf}$ gives $S$, $g$ and
$p_0$, and the position is a Gaussian of width $\sqrt{2\sigma}/\pi$ about $\xi$ with the
two nearest images removed by thinning; above $\sigma_{\max}$ only the slowest mode
survives, $\ln S$ falls linearly in $\sigma$, $g=1$, $p_0=1/2$, and the $\sigma_{\max}$
positions are reused.  The box survival at the budget, $\ln S_b=\sum_i\ln S_i$, and the
inversion for the exit time are done on the fly by bisection in $\ln t$, sixty
iterations of three table lookups each.

\subsection{The sphere, for the record}

The sphere branch tabulates the first-passage CDF $F(s)=1-P(s)$,
$P(s)=2\sum_n(-1)^{n+1}e^{-n^2s}$, on $\ln s\in[\ln10^{-3},\ln40]$, and the radius of a
walker still inside, $p(\rho\,|\,s)\propto\rho\sum_nn\sin(n\pi\rho)e^{-n^2s}$, as
quantiles on $(\ln s,\text{level})$ from $s=0.03$ up; below that the free-space Gaussian
in $r^2=2Dct\,\chi^2_3$ is exact to $e^{-\pi^2/4s}$.

\section{Energy, weight and the scatterings the step stands for}
\label{sec:energy}

With coherent scattering (Thomson, isotropic) the energy is unchanged and the weight is
multiplied by $\exp(-\alpha_{\rm abs}ct)$.  With Compton scattering the step uses the
quantile table named by \key{kgreens\_file}: $y=\theta\,n_e\sigma_Tct$ with the Thomson
coefficient of the cell, $x_i=E/kT$ the comoving energy, $\Lambda$ the free-free
absorption parameter built from the same per-cell prefactor the opacity uses, and the
table gives $x_f$ at a uniform level and $S$.  The weight is multiplied by $S$; the
ordinary \key{minweight} roulette applies afterwards as it does everywhere.  Where the
absorption is not free-free the $\Lambda=0$ slice is used together with the grey factor.
The scattering count the path implies, $\chi_{\rm sc}ct$, is added to the photon's and
to the block's \texttt{nmrw\_scat}, which the run report prints beside the steps taken.

The step declines, and the photon scatters normally, outside the Kompaneets description:
$\theta>0.03$, $x_i\theta>0.05$, $x_i$ outside the table, or
$\theta\Lambda a(x_i)>0.1$ (absorption too strong for energy and space to separate).  The
first two are also tested in the trigger, before any frame is built, since a hot corona
fails them at every interaction and the frame is the expensive part under the general
pusher.

\section{Moments}
\label{sec:moments}

The step contributes to the lab moments of the cell it walked in as an isotropic
comoving radiation field with energy density estimate
\begin{equation}
  E_c=\frac{w\,\bar E\,ct}{c},\qquad \bar E=\tfrac12(E_{\rm before}+E_{\rm after}),
\end{equation}
the path-length estimator with the whole path, at the mean of the two energies.  Seen
from the lab frame in which the fluid moves at $\boldsymbol\beta$ with Lorentz factor
$\gamma$,
\begin{equation}
  E=\gamma^2\Bigl(1+\tfrac{\beta^2}{3}\Bigr)E_c,\qquad
  F^i=\tfrac43\gamma^2\beta^i cE_c,\qquad
  P^{ij}=E_c\Bigl(\tfrac13\delta^{ij}+\tfrac43\gamma^2\beta^i\beta^j\Bigr),
\end{equation}
which \texttt{MonteCarloBlock::UpdateMomentsMRW} deposits into \texttt{MCIER},
\texttt{MCIFR}$_i$ and the six pressure components.  Under the general pusher the lab
frame is the normal observer's and $\boldsymbol\beta$ is the fluid's velocity on its
tetrad legs.  This is exact in expectation over many steps and approximate for one; the
scattering-source moments (\texttt{mom\_flag\_scat}), which bin the energy, receive
nothing from the step, so the Compton source terms of a cell that is mostly walked are
under-sampled.  In the static post-processing runs the step was built for, the moments
are not the product; in a coupled run that would need attention.

\section{A moving medium: the comoving frame and the advection room}
\label{sec:advection}

The walk is done in the comoving frame, where the medium is at rest and the Green's
function applies.  The photon enters it with the block loop's own
\texttt{TransformToComoving} and leaves it with \texttt{TransformToCoordinate}, so its
energy, direction and opacities are in the right frame without new code; what is new is
the geometry.

Over comoving time $t$ the fluid element moves $\boldsymbol\beta\gamma ct$ in the lab,
and the cell is fixed in the lab, so the domain plus that displacement has to stay inside
the cell.  A fraction $f$ of the nearest-face distance $d$ is set aside for advection
and the rest is the domain.  $f$ is chosen so that the expected diffusion time across
the domain, $(1-f)^2d^2\chi/2$, carries the centre just the distance set aside,
$fd=\beta\gamma c\,t$, i.e.\ with $a=\tfrac12\beta\gamma\chi d$,
\begin{equation}
  f=\frac{(2a+1)-\sqrt{4a+1}}{2a},\qquad f\le\tfrac12,
\end{equation}
small for a slow flow and approaching the cap when advection would cross the cell before
the walk could leave even a small domain, in which case the thickness test fails and the
step declines (counted as ``too thin'').  The box is shrunk on its downstream side in
each dimension by $fd\,|\beta_i|/|\boldsymbol\beta|$ and the walk takes place in the
shrunken box; the comoving time budget is $t_b=fd/(\beta\gamma c)$, or less if the cycle
remainder is shorter.  On exit the lab displacement along the box axes is the comoving
one plus $\boldsymbol\beta\gamma ct_e$, and the lab time advances by $\gamma ct_e$.
Lorentz contraction of the box is neglected.

On the gridded sphere with a uniform flow of $0.1c$ and $0.3c$ through the fixed escape
surface, the walk agreed with the analog run of the same deck to $0.7$ and $0.6$ of the
two-sample Kolmogorov scale in the escaping spectrum, with mean escaping energies
1.2 and 0.6 percent low and escape paths 1.5 and 2 percent short; the energy deficit
follows the path shortening (Section~\ref{sec:thresholds}) rather than growing with
$\beta$.

\section{Cartesian cells, and spherical-polar cells}
\label{sec:cells}

In a Cartesian cell the box is the cell and the box axes are the coordinate axes; the
widths and the photon's offsets from the lower faces come straight from the face
coordinates (\texttt{PhotonPusher::CellGeometry}), in code length, and the displacement
is added to the position.

A spherical-polar cell is treated as the box
\begin{equation}
  W_r=\Delta r,\qquad W_\theta=r\,\Delta\theta,\qquad W_\phi=r\sin\theta\,\Delta\phi
\end{equation}
at the photon's $r$ and $\theta$, with offsets $r-r_f$, $r(\theta-\theta_f)$,
$r\sin\theta(\phi-\phi_f)$, in the pusher's local orthonormal $(\hat r,\hat\theta,\hat\phi)$
basis, the one its directions and the per-cell velocity \texttt{vel} are expressed in.
The walk's displacement, with the advection, is then a vector in that basis; it is
carried to Cartesian components at the photon,
\begin{equation}
  \begin{pmatrix}e_x\\e_y\\e_z\end{pmatrix}=
  \begin{pmatrix}\sin\theta\cos\phi&\cos\theta\cos\phi&-\sin\phi\\
                 \sin\theta\sin\phi&\cos\theta\sin\phi&\cos\phi\\
                 \cos\theta&-\sin\theta&0\end{pmatrix}
  \begin{pmatrix}d_r\\d_\theta\\d_\phi\end{pmatrix},
\end{equation}
added to the photon's Cartesian position, and converted back to $(r,\theta,\phi)$.  The
legacy implementation did this with sign errors in the $\hat\theta$ and $\hat\phi$
columns; the matrix above is the orthonormal basis of spherical coordinates written out.

What the box approximation means for a curvilinear cell: the faces are not planes.  The
$\theta$ and $\phi$ faces of a cell are surfaces of constant angle, and a straight
displacement in the local basis of the photon's starting point ends at a different $r$
and $\theta$ from the one the box arithmetic assumed, by amounts of order the cell's
angular extent times its size.  For cells with $\Delta\theta$ and $\Delta\phi$ of a few
degrees, and away from the axis and the origin, this is a few percent of the cell
width, the same order as the geodesic curvature across a Kerr-Schild cell that the
general pusher's step also neglects, and it is caught in the same way: the cell update
after the step finds the photon's true cell, and the boundary functions run if the
rounding or the curvature put it across a block face.  Where the approximation is poor
is where the cell itself is poorly described as a box, within a few cells of the origin
in $r$, where $W_\theta$ varies across the cell by the ratio of the inner to the outer
radius, and near the polar axis, where $W_\phi\to0$ and the general pusher has defects
of its own that the polar-axis notes record.  The trigger will seldom fire there, since
such cells are thin in at least one direction, but a problem that is thick near its
origin or axis should mask those cells (\texttt{accel\_mask}) rather than trust the box.
The spherical-polar branch of the step has run through the pusher's own tests and the
sphere problem on a Cartesian grid; a Thomson atmosphere in spherical polar coordinates
with boosts on, the \texttt{mc\_isoth} deck of the velocity reconstruction plan, is the
regression that exercises it directly and is still to be run.

\section{The general pusher: the step in curved spacetime}
\label{sec:gr}

Under the general pusher the lab frame is the normal observer's orthonormal tetrad at the
photon, built from the metric pair the pusher caches at the end of its last step
(\texttt{MetricPairAt}, \texttt{NormalObserver}, \texttt{ConstructTetrad}).  The cell's
three coordinate edges, $(0,\Delta x^i,0,0)$ and permutations, are projected into that
tetrad; their lengths are the box widths $W_i$ and the photon's offsets are its
coordinate offsets scaled by $W_i/\Delta x^i$.  The fluid's four-velocity at the photon
(\texttt{FluidFourVelocity}; in flat spacetime the per-cell \texttt{vel} carried onto
coordinate components with the static tetrad, as \texttt{ComovingFrame} does) gives
$\gamma=u^{\hat0}$ and the velocity on the tetrad's spatial legs, which is what the
moments use.  The drift that sets the advection room is a different quantity: the cell
is fixed in coordinates, so the fluid's drift relative to it is its coordinate velocity
per unit lab proper time, $u^i/\gamma$, which carries the shift as well as the boost, and
along box axis $i$ it is $u^i/\gamma$ in units of $W_i/\Delta x^i$.  The step then runs
exactly as in Section~\ref{sec:advection} with that $\boldsymbol\beta$.

On exit the displacement along each box axis is a fraction of the corresponding
coordinate edge, so the coordinate position advances by $db_i\,\Delta x^i/W_i$, and the
coordinate time by the lab proper time over the lapse, $\gamma ct_e/\alpha$.  The new
direction is drawn isotropically in the comoving frame and carried to coordinate
components by \texttt{TransformToCoordinate} with the comoving tetrad \emph{at the exit
point}, which is where the change of redshift and of fluid velocity across the step
enters, to first order in the cell size; nothing in the step itself knows about the
metric's variation.  Curvature within the cell is neglected.  A polarized photon comes
out of the step unpolarized: its Stokes intensity is read off the coherency tensor before
the walk, $Q=U=V=0$ afterwards, and the tensor rebuilt for the new direction before the
transform back.

Two things learned the hard way, both fixed: the transform back to coordinate components
must happen before the cell update, since the update can hand the photon to another
block, which expects coordinate components, and a photon that crossed carrying a
comoving unit direction came out with energies of either sign and time running
backwards; and the advection room has to use the full coordinate drift $u^i/\gamma$,
not the velocity relative to the normal observer alone, since near the hole the shift is
as large as the boost.

In flat spacetime with \key{general\_pusher = true} the same code runs with the
Minkowski tetrad, and on the gridded sphere the general pusher's walk agreed with its
analog run to $0.6$ (spectrum) and $0.8$ (paths) of the two-sample scale at rest and
$0.5$ and $2.5$ at $0.1c$; the two pushers agree with each other on the same decks.  On
the high-resolution X-ray binary snapshot in Kerr-Schild coordinates the accelerated run
reproduced the emission run's bands within the output-to-output scatter
($0.97\pm0.02$ at 0.1 to 0.3 keV, $0.98\pm0.04$, $1.02\pm0.03$, $1.04\pm0.11$ at 3 to
10 keV) with 6 percent of the scatterings replaced.

\section{After the step: the step never decides the cell}
\label{sec:after}

The step moves the photon; it does not decide which cell the photon ends in.  That is
left to the same cell update (\texttt{UpdateZone}) that follows any flight, and the order
of operations is fixed: the moments are deposited in the cell the walk happened in; the
photon is carried back from the comoving frame to the lab with the frame of that cell,
at its new position; only then is its cell found.

Under the general pusher the box is the cell by construction: each displacement is a
fraction of the box width applied to the cell's coordinate width, so a walk that ends on
a wall ends on the face (plus a margin small enough that the update sees the crossing)
and a walk that ends inside ends inside.  What is approximate is the shape of the domain
the photon diffused in (edges measured in the tetrad at the starting point and treated
as orthogonal, the metric and fluid velocity constant over the cell), not where it is
put relative to the faces.  Under the legacy spherical-polar pusher the box's angular
walls are flat while the faces are surfaces of constant angle, so a walk that ends on a
wall can land slightly inside or slightly beyond the face.

In every case the outcome is decided by where the photon actually is.  If it is still in
the walked cell it continues there as an ordinary photon with a fresh optical depth;
being within the face threshold of a wall, it scatters normally rather than stepping
again until it has wandered inward.  If it is across a face, its indices move to the
neighbour, and at a block face the boundary function buffers it to the other rank or
ends it, exactly as for a photon whose flight crossed the face.  The legacy pushers then
recompute the opacities in the new cell at the lab energy, as at any crossing, and the
general pusher's loop refreshes its opacities through the comoving shift on its next
step and discards its cached metric and connection.  A photon that lands in a
neighbour whose fluid velocity differs was carried back to the lab with the walked
cell's velocity; that is the same inconsistency any photon crossing a face mid-flight
has, of order the velocity jump between neighbours.

\section{Triggers, thresholds and declines}
\label{sec:thresholds}

The trigger (\texttt{MRWTrigger}) is evaluated at every interaction that would fall
inside the cell and must be cheap.  With the cell domain it requires the cell's optical
half-width, in the lab extinction at the photon's energy, to exceed \key{accel\_tau}
(default 20) and the photon's distance to every face to exceed \key{accel\_face\_tau}
(default 5); with the sphere domain, the face distance to exceed \key{accel\_tau}.  A
cell masked by the problem generator (\texttt{accel\_mask}) never triggers; the sphere
test masks the cells its escape surface cuts, which are not the uniform medium the step
assumes.  The Compton pre-checks of Section~\ref{sec:energy} follow.

The step itself may still decline once the comoving extinction and the advection room
are known (``too thin''), when there is no usable frame (inside a horizon), when the
Kompaneets guards that need $\Lambda$ fail, or when the time budget is exhausted; the
report counts each.  Under the general pusher an attempt costs a tetrad and two
transforms, so on a grid where most attempts decline the step can cost more than it
saves; on the X-ray binary 69 percent of attempts declined and the run was 5 percent
slower than the analog one.

The face threshold is where the diffusion description is weakest.  The first-passage
time of diffusion theory lets a photon leave arbitrarily soon, whereas a real photon two
mean free paths from a face needs at least that path to reach it; steps taken close to a
face therefore come out short.  On the gridded sphere at $\tau=100$ the escape paths
were 3 to 4 percent short with the threshold at 2 (at $2.7$ times the two-sample scale)
and consistent with the analog run at 5, hence the default; at the same time the step
replaced 61 and 45 percent of the scatterings.  In a Compton-thick medium the
shortening shows up as a lower mean escaping energy, since $y$ is proportional to the
path, which is what the moving-medium tests saw at the percent level.  A scan of the
threshold at a Compton-thick setting is the measurement that would replace the default
with a number for production.

\section{Keys and report lines}
\label{sec:keys}

\begin{center}
\begin{tabular}{llp{8.2cm}}
\toprule
key (\texttt{<montecarlo>}) & default & meaning \\
\midrule
\key{acceleration} & false & turn the step on \\
\key{accel\_domain} & \texttt{cell} & \texttt{cell} or \texttt{sphere} \\
\key{accel\_tau} & 20 & cell optical half-width (cell domain) or face distance (sphere) a step needs \\
\key{accel\_face\_tau} & 5 & cell domain: a photon closer than this to a face scatters normally \\
\key{accel\_pmax} & 0 & cap the time budget where the domain survival has fallen to this; 0 is no cap \\
\key{kgreens\_file} & -- & the Kompaneets quantile table, a flat binary from \texttt{kgreens\_export.py}; required with Compton scattering \\
\key{accel\_report} & false & histogram the analog scatterings by the optical half-width of their cell \\
\bottomrule
\end{tabular}
\end{center}

The report prints, per transport, the steps taken, the scatterings their paths stand
for, and the declines by reason; with \key{accel\_report}, the histogram that says how
much of a run's work a step confined to one cell can reach.  On the X-ray binary that
histogram put 9 percent of the scatterings in cells thick enough to step in at the
default threshold and 60 percent in cells of one to sixteen mean free paths, which is
why the in-cell step is a correctness result there and not a saving; the plan
(\texttt{doc/monte\_carlo/compton\_mrw\_plan.md}, Section~6.7) records the region step
that would reach those cells and why it waits.

\section{Verification summary}

\begin{center}
\small
\begin{tabular}{llrll}
\toprule
case & domain & replaced & spectrum & paths \\
\midrule
sphere $\tau=100$, $8^3$ cells, at rest & sphere & 13\% & agrees & 1\% short \\
same, cell domain, face threshold 5 & cell & 45\% & 0.7 & 1.1 \\
sphere $\tau=1000$, cell domain & cell & 81\% & 0.9 & 1.2 \\
sphere $\tau=100$, $0.1c$ flow & cell & 35\% & 0.7 & 1.6 \\
general pusher, flat, at rest & cell & 44\% & 0.6 & 0.8 \\
general pusher, flat, $0.1c$ & cell & 35\% & 0.5 & 2.5 \\
Thomson sphere $\tau=100$ & cell & 54\% & -- & 0.9 \\
X-ray binary, Kerr-Schild & cell & 6\% & bands within scatter & -- \\
\bottomrule
\end{tabular}
\end{center}
The spectrum and path columns are the maximum CDF difference in units of the two-sample
Kolmogorov scale, $\sqrt{(n_a+n_b)/n_an_b}$, between the walk and the analog run of the
same deck; the X-ray binary comparison is band luminosities against the emission run's
output-to-output scatter.  Every case also passes the Green's function tests of the
sphere problem (\texttt{tst/montecarlo/sphere\_compton}) where those apply, and the C++
table sampler reproduces the Python reference to $10^{-14}$.

\end{document}
